% Generated by qpaper-workflow from paper/source/questions.yaml for core4

\section{Introduction}
% <qpaper-block id="para.problem.aac_text_entry" question="Q.problem.aac_text_entry" claims="" evidence="">
Hands-free text entry is a central problem for augmentative and alternative communication (AAC): many users and usage settings cannot assume reliable speech, hand movement, or a visible keyboard. Silent speech interfaces offer an appealing route, but a practical AAC channel must support more than a small set of pre-defined commands. It must let users compose words and messages.

% </qpaper-block>

% <qpaper-block id="para.problem.epg_signal" question="Q.problem.epg_signal" claims="C.epg_signal" evidence="A_results_paper_layout_2026_03_03_table_dataset_summary_csv,A_results_dataset_audit_silentspeller_2026_02_24_dataset_summary_csv,A_results_dataset_audit_silentspeller_2026_02_24_report_md">
Electropalatography (EPG) makes this setting concrete. A SmartPalate-style sensor records tongue-palate contact as a multichannel temporal pattern while the user silently articulates. In the SilentSpeller-style recordings used in this work, each trial is a variable-length EPG sequence paired with the word that the participant attempted to spell or articulate.

% </qpaper-block>

% <qpaper-block id="para.problem.prior_capability" question="Q.problem.prior_capability" claims="" evidence="A_results_related_work_survey_epgspeller_2026_06_27_v2_report_md,A_results_related_work_survey_epgspeller_2026_06_27_v2_coverage_matrix_csv,A_results_related_work_survey_epgspeller_2026_06_27_v2_references_audit_md">
Prior silent-speech interfaces have shown that silent articulatory or biosignal patterns can be mapped to useful commands or constrained text-entry actions. That capability is important, but it is not the same as an AAC text channel that lets a user spell arbitrary messages from character evidence.

% </qpaper-block>

% <qpaper-block id="para.problem.fixed_vocabulary" question="Q.problem.fixed_vocabulary" claims="C.fixed_vocab_gap" evidence="A_results_related_work_survey_epgspeller_2026_06_27_v2_report_md,A_results_related_work_survey_epgspeller_2026_06_27_v2_coverage_matrix_csv,A_results_related_work_survey_epgspeller_2026_06_27_v2_references_audit_md">
The key interface question is therefore not simply whether EPG can recognize a known silent command. EPGSpeller asks whether the contact signal can be decoded at the character level. Character-level silent spelling breaks the fixed-vocabulary limitation of command recognition: if the spelling stream is recognized well enough, the user is no longer restricted to words that were present as output classes during training. The related-work comparison also motivates keeping command recognition, modality-specific SSI systems, and LM-assisted text entry separate from the present lexicon-free EPG spelling claim.

% </qpaper-block>

% <qpaper-block id="para.claim.central" question="Q.claim.central" claims="C.central" evidence="A_results_paper_layout_2026_03_03_table_main_compact_csv,A_results_msx20260224_paper_tables_table_main_csv,A_results_slt_p2_20260626_paper_tables_table_silentspeller_lite_csv,A_results_slt_p2_20260626_paper_tables_table_kshot_curve_csv,A_results_slt_p2_20260626_kshot_feasibility_csv,A_results_preproc20260626_paper_tables_table_preproc_protocol_condition_csv,A_results_preproc20260626_preproc_ablation_report_md">
This paper presents EPGSpeller, a character-level recognizer for electropalatographic silent spelling. The contribution is not only a lower-error decoder. The system shows that EPG can be framed as a hands-free AAC text-entry interface that is no longer tied to a fixed command vocabulary. At the same time, the experiments deliberately expose the remaining bottleneck: personalized calibration is effective, whereas cross-participant transfer and high-shot calibration are still constrained by the available data. The preprocessing sweep reinforces this interpretation by showing that simple normalization changes cross-user errors but does not turn them into personalized recognition.

% </qpaper-block>

% <qpaper-block id="para.claim.evidence_preview" question="Q.claim.evidence_preview" claims="" evidence="A_results_paper_layout_2026_03_03_table_main_compact_csv,A_results_slt_p2_20260626_paper_tables_table_silentspeller_lite_csv,A_results_slt_p2_20260626_paper_tables_table_kshot_curve_csv,A_results_preproc20260626_paper_tables_table_preproc_protocol_condition_csv">
The evidence is organized around that claim boundary. The paper first tests personalized and word-holdout spelling recognition, then asks whether calibration can be reduced, and finally checks whether multi-source transfer or simple preprocessing removes the user-dependence.

% </qpaper-block>

% <qpaper-block id="para.claim.contributions" question="Q.claim.contributions" claims="" evidence="A_results_paper_layout_2026_03_03_table_main_compact_csv,A_results_slt_p2_20260626_paper_tables_table_silentspeller_lite_csv,A_results_slt_p2_20260626_kshot_feasibility_csv">
The contributions are: (1) an AAC-centered framing of SmartPalate EPG as a character-level silent spelling channel, (2) a high-performing lexicon-free vector/GRU recognizer for that task, and (3) evidence that connects personalized spelling, SilentSpeller-compatible word holdout, and calibration limits without hiding the remaining user-dependence behind a single benchmark score.

% </qpaper-block>

\section{Task and Method}
% <qpaper-block id="para.method.signal_representation" question="Q.method.00_signal_representation" claims="" evidence="A_results_paper_layout_2026_03_03_table_dataset_summary_csv,A_results_dataset_audit_silentspeller_2026_02_24_dataset_summary_csv">
The input representation is the EPG contact stream itself: each trial is treated as an ordered sequence of palate-contact frames rather than as audio, video, or a manually segmented phone sequence. This keeps the task tied to the SmartPalate-style signal and makes the decoding claim about silent tongue-palate contact.

% </qpaper-block>

% <qpaper-block id="para.task.definition" question="Q.method.01_task_definition" claims="C.task" evidence="A_results_paper_layout_2026_03_03_table_dataset_summary_csv,A_results_msx20260224_paper_tables_table_main_csv,A_results_slt_p2_20260626_paper_tables_table_silentspeller_lite_csv">
Given an EPG contact sequence, the model outputs a character string. We treat greedy character-level decoding as the primary lexicon-free result. Lexicon-projected scores are useful diagnostics, but they are not allowed to carry the central claim because an AAC spelling interface must eventually support words and messages outside a fixed evaluation lexicon.

% </qpaper-block>

% <qpaper-block id="para.method.workflow" question="Q.method.02_workflow" claims="C.workflow" evidence="A_results_msx20260224_paper_tables_table_main_csv,A_results_msx20260224_paper_tables_table_spatial_k124_csv">
EPGSpeller uses a character-level CTC recognizer over EPG frames. We evaluate a high-capacity vector/GRU recognizer as the main system and compare it with layout-aware variants that use row-column pooling or two-dimensional palatal grids. The design goal is deliberately pragmatic: maximize silent spelling accuracy under lexicon-free decoding, then use ablations to ask which EPG-specific inductive biases help under the same protocols.

% </qpaper-block>

% <qpaper-block id="para.method.variant_scope" question="Q.method.03_variant_scope" claims="" evidence="A_results_msx20260224_paper_tables_table_spatial_k124_csv,A_results_msx20260224_paper_tables_table_k64_methods_csv">
The method scope is therefore intentionally narrow. The vector/GRU system is the main recognizer used for the paper's interface claim, while layout-aware and reduced-electrode variants are treated as controlled engineering ablations under the same decoding and reporting rules.

% </qpaper-block>

\section{Evaluation}
% <qpaper-block id="para.setup.data" question="Q.setup.data_protocols" claims="C.protocols" evidence="A_results_paper_layout_2026_03_03_table_dataset_summary_csv,A_results_dataset_audit_silentspeller_2026_02_24_dataset_summary_csv,A_results_msx20260224_paper_tables_table_main_csv,A_results_slt_p2_20260626_paper_tables_table_lowshot_extended_csv,A_results_slt_p2_20260626_paper_tables_table_kshot_curve_csv,A_results_slt_p2_20260626_paper_tables_table_silentspeller_lite_csv">
The evaluation uses four SilentSpeller-style EPG datasets with different repetition densities: p1 and p2 each contain 2328 trials over 1164 word types, p3 contains 2797 trials over 1164 word types, and p4 contains 1167 trials over 1162 word types. The dataset summary table supports the feasibility decisions used later in the calibration analysis, because low-shot and high-shot calibration depend on repeated productions per word, not only on the number of participants.

% </qpaper-block>

% <qpaper-block id="tab.dataset_summary" question="Q.setup.data_protocols" claims="C.protocols" evidence="A_results_paper_layout_2026_03_03_table_dataset_summary_csv">
\begin{table}[t]
\centering
\caption{Dataset audit summary for the SilentSpeller-style EPG recordings.}
\label{tab:dataset_summary}
\begin{tabular}{llllll}
\hline
id & N & V & Tmed & contact & zero \\
\hline
p1 & 2328 & 1164 & 234 & 0.19 & 0 \\
p2 & 2328 & 1164 & 196 & 0.20 & 0 \\
p3 & 2797 & 1164 & 283 & 0.12 & 4 \\
p4 & 1167 & 1162 & 229 & 0.19 & 0 \\
\hline
\end{tabular}
\end{table}
% </qpaper-block>

% <qpaper-block id="para.setup.protocols" question="Q.setup.protocols" claims="" evidence="A_results_msx20260224_paper_tables_table_main_csv,A_results_slt_p2_20260626_paper_tables_table_lowshot_extended_csv,A_results_slt_p2_20260626_paper_tables_table_kshot_curve_csv,A_results_slt_p2_20260626_paper_tables_table_silentspeller_lite_csv">
The evaluation separates the claims that matter for an AAC spelling interface. Personalized protocols measure whether the recognizer can work after user-specific training; word-holdout protocols ask whether spelling generalizes beyond training word types; cross-participant protocols test whether a model can be reused without target calibration; and k-shot protocols measure how much target-user calibration is needed. Reported tables aggregate repeated runs by condition and keep incompatible protocols separate rather than pooling instance holdout, word holdout, and cross-participant transfer into one score.

% </qpaper-block>

% <qpaper-block id="para.setup.metrics" question="Q.setup.metrics" claims="C.metrics" evidence="A_results_paper_layout_2026_03_03_table_main_compact_csv,A_results_msx20260224_paper_tables_table_main_csv,A_results_slt_p2_20260626_paper_tables_table_silentspeller_lite_csv">
The primary metric is greedy character error rate (CER), computed without projecting the output onto a candidate lexicon. Where useful, we also report lexicon-projected scores to compare with closed-vocabulary behavior, but those values are not used as the main evidence for lexicon-free text entry. Real-time factor (RTF) is reported separately to show that decoding latency is not the main obstacle in the present experiments.

% </qpaper-block>

% <qpaper-block id="para.setup.baselines_controls" question="Q.setup.controls" claims="" evidence="A_results_paper_layout_2026_03_03_table_main_compact_csv,A_results_slt_p2_20260626_paper_tables_table_silentspeller_lite_csv,A_results_preproc20260626_paper_tables_table_preproc_protocol_condition_csv,A_results_msx20260224_paper_tables_table_spatial_k124_csv,A_results_msx20260224_paper_tables_table_k64_methods_csv">
Baselines and controls are separated by what they test. The main vector/GRU path carries the recognition result; SilentSpeller-compatible-lite word holdout connects to the historical task; lexicon projection is reported only as a closed-vocabulary diagnostic; and transfer and preprocessing controls are used to bound the calibration claim rather than to redefine the primary task.

% </qpaper-block>

\section{Results}
% <qpaper-block id="para.result.personalized" question="Q.result.01_personalized" claims="C.personalized_result" evidence="A_results_paper_layout_2026_03_03_table_main_compact_csv,A_results_msx20260224_paper_tables_table_main_csv">
Table~\ref{tab:main_protocols} gives the central recognition result. With the same vector/GRU recognizer and greedy decoding, personalized conditions are much stronger than cross-participant transfer: P1 reaches 0.1796 CER, P2 reaches 0.1451 CER, whereas P3 cross-participant transfer remains at 0.6909 CER. The contrast is the main empirical shape of the paper: EPG silent spelling is viable after user-specific calibration, but it is not yet a calibration-free interface.

% </qpaper-block>

% <qpaper-block id="tab.main_protocols" question="Q.result.01_personalized" claims="C.personalized_result" evidence="A_results_paper_layout_2026_03_03_table_main_compact_csv">
\begin{table}[t]
\centering
\caption{Main protocol recap for the vector/GRU recognizer. CER is greedy character error rate; lex_all is lexicon-projected error.}
\label{tab:main_protocols}
\begin{tabular}{llll}
\hline
protocol & n & cer & lex \\
\hline
P1 & 4 & 0.180±0.084 & 0.102±0.068 \\
P2 & 3 & 0.145±0.063 & 0.075±0.048 \\
P3 & 6 & 0.691±0.133 & 0.644±0.060 \\
\hline
\end{tabular}
\end{table}
% </qpaper-block>

% <qpaper-block id="para.result.silentspeller_lite" question="Q.result.02_silentspeller_lite" claims="C.silentspeller_lite" evidence="A_results_slt_p2_20260626_paper_tables_table_silentspeller_lite_csv,A_results_slt_p2_20260626_slt_p2_report_md">
To connect this framing back to SilentSpeller-style evaluation, we additionally use a word-holdout lite condition on subj1 with 100 unseen test word types. The vector/GRU recognizer obtains 0.1133 +/- 0.0128 greedy CER across four seeds. This is not claimed as a full recreation of the historical implementation; it is a same-dataset, spelling-oriented bridge showing that the recognizer remains effective when the test words are held out.

% </qpaper-block>

% <qpaper-block id="tab.silentspeller_lite" question="Q.result.02_silentspeller_lite" claims="C.silentspeller_lite" evidence="A_results_slt_p2_20260626_paper_tables_table_silentspeller_lite_csv">
\begin{table}[t]
\centering
\caption{SilentSpeller-compatible-lite word-holdout result for subj1.}
\label{tab:silentspeller_lite}
\begin{tabular}{lllllll}
\hline
condition & n_test_words & n & cer_m & cer_s & rtf_m & rtf_s \\
\hline
subj1_test100 & 100 & 4 & 0.113343 & 0.012833 & 0.000290 & 0.000042 \\
\hline
\end{tabular}
\end{table}
% </qpaper-block>

% <qpaper-block id="para.result.calibration" question="Q.result.03_calibration" claims="C.calibration_curve" evidence="A_results_slt_p2_20260626_paper_tables_table_kshot_curve_csv,A_results_slt_p2_20260626_paper_tables_table_lowshot_extended_csv,A_results_slt_p2_20260626_kshot_feasibility_csv">
The calibration experiments turn the cross-user weakness into a more actionable result. For subj3, zero-shot transfer gives 0.8403 CER for single-source transfer and 0.8381 CER for multi-source transfer, while one and two target examples per word reduce CER to 0.3503 and 0.2458. Additional k=1 runs on subj1 and subj2 give 0.1082 and 0.1118 CER, indicating that one-shot personalized calibration can be strong when each word has two productions. However, the available public data do not support a standard 50-word held-out evaluation at k=4 or k=8: only 14 word types satisfy the k=4 repetition requirement for subj3, and none satisfy k=8.

% </qpaper-block>

% <qpaper-block id="tab.kshot_curve" question="Q.result.03_calibration" claims="C.calibration_curve" evidence="A_results_slt_p2_20260626_paper_tables_table_kshot_curve_csv">
\begin{table}[t]
\centering
\caption{Calibration curve and feasibility for subj3. The k=4 row is reduced-vocabulary diagnostic evidence; k=8 is data-infeasible.}
\label{tab:kshot_curve}
\begin{tabular}{lllllllll}
\hline
target & k & condition & n & cer_m & cer_s & rtf_m & rtf_s & notes \\
\hline
subj3 & 0 & P3_single_source & 8 & 0.840301 & 0.094775 & 0.000097 & 0.000004 & msx/uc vector \\
subj3 & 0 & P3MS_multi_source & 4 & 0.838062 & 0.033612 & 0.000096 & 0.000002 & msx:p3ms vector \\
subj3 & 1 & P2K_within_subject & 4 & 0.350280 & 0.059950 & 0.000164 & 0.000009 & msx:existing \\
subj3 & 2 & P2K_within_subject & 4 & 0.245802 & 0.008507 & 0.000165 & 0.000002 & msx:existing \\
subj3 & 4 & P2K_reduced_vocab & 4 & 0.781609 & 0.208381 & 0.002482 & 0.000169 & sltp2:new reduced vocabulary (min_repetitions=5, n_competition_words=1) \\
subj3 & 8 & N/A & 0 &  &  &  &  & infeasible (feasible_vocab=0, feasible_for_50_word_eval=False) \\
\hline
\end{tabular}
\end{table}
% </qpaper-block>

% <qpaper-block id="tab.lowshot_extended" question="Q.result.03_calibration" claims="C.calibration_curve" evidence="A_results_slt_p2_20260626_paper_tables_table_lowshot_extended_csv">
\begin{table}[t]
\centering
\caption{Low-shot extension beyond subj3 for subjects with two repetitions per word.}
\label{tab:lowshot_extended}
\begin{tabular}{llllllll}
\hline
subject & k & n & cer_m & cer_s & rtf_m & rtf_s & source \\
\hline
subj1 & 1 & 4 & 0.108222 & 0.005059 & 0.000106 & 0.000001 & sltp2:new \\
subj2 & 1 & 4 & 0.111769 & 0.006698 & 0.000117 & 0.000001 & sltp2:new \\
subj3 & 1 & 4 & 0.350280 & 0.059950 & 0.000164 & 0.000009 & msx:existing \\
subj3 & 2 & 4 & 0.245802 & 0.008507 & 0.000165 & 0.000002 & msx:existing \\
\hline
\end{tabular}
\end{table}
% </qpaper-block>

% <qpaper-block id="para.result.cross_user" question="Q.result.04_cross_user" claims="C.cross_user_bottleneck" evidence="A_results_paper_layout_2026_03_03_table_p3ms_compact_csv,A_results_paper_layout_2026_03_03_table_to4_compact_csv,A_results_msx20260224_paper_tables_table_p3ms_csv,A_results_uc20260226_paper_tables_table_to4_generalization_csv,A_results_preproc20260626_paper_tables_table_preproc_protocol_condition_csv,A_results_preproc20260626_preproc_ablation_report_md">
Multi-source transfer does not remove the bottleneck. Paired-source and multi-source variants sometimes improve over a single-source baseline, but the resulting errors are still far from personalized performance and vary strongly by target. For target subject 4, all-to-4 vector CER is 0.686 in P3 and 0.582 in P3MS. A preprocessing sweep over raw, standardized, contact-matched, target-marginal, and PCA-standardized variants improves the best protocol-level averages only to 0.6426 for P3 and 0.5692 for P3MS. For the paper's AAC framing, this means the current evidence supports personalized silent spelling, not an off-the-shelf model that can be deployed to a new user without calibration.

% </qpaper-block>

% <qpaper-block id="tab.p3ms" question="Q.result.04_cross_user" claims="C.cross_user_bottleneck" evidence="A_results_paper_layout_2026_03_03_table_p3ms_compact_csv">
\begin{table}[t]
\centering
\caption{Multi-source cross-participant transfer results.}
\label{tab:p3ms}
\begin{tabular}{llll}
\hline
group & variant & cer & rtf \\
\hline
p23->p1 & vec & 0.473±0.014 & 0.0001±0.0000 \\
p23->p1 & rowcol & 0.476±0.014 & 0.0002±0.0000 \\
p23->p1 & grid & 0.602±0.076 & 0.0004±0.0000 \\
p23->p1 & grid_aug & 0.691±0.165 & 0.0004±0.0000 \\
p13->p2 & vec & 0.520±0.003 & 0.0001±0.0000 \\
p13->p2 & rowcol & 0.504±0.012 & 0.0003±0.0000 \\
p13->p2 & grid & 0.511±0.047 & 0.0004±0.0000 \\
p13->p2 & grid_aug & 0.511±0.016 & 0.0004±0.0000 \\
p12->p3 & vec & 0.838±0.034 & 0.0001±0.0000 \\
p12->p3 & rowcol & 0.790±0.033 & 0.0002±0.0000 \\
p12->p3 & grid & 0.736±0.018 & 0.0003±0.0000 \\
p12->p3 & grid_aug & 0.756±0.020 & 0.0003±0.0000 \\
\hline
\end{tabular}
\end{table}
% </qpaper-block>

% <qpaper-block id="tab.to4_generalization" question="Q.result.04_cross_user" claims="C.cross_user_bottleneck" evidence="A_results_paper_layout_2026_03_03_table_to4_compact_csv">
\begin{table}[t]
\centering
\caption{Additional cross-participant generalization results for target subject 4.}
\label{tab:to4_generalization}
\begin{tabular}{lllll}
\hline
proto & lvl & group & cer & lex \\
\hline
P3 & dir & p1->p4 & 0.666±0.024 & 0.651±0.026 \\
P3 & dir & p2->p4 & 0.809±0.021 & 0.736±0.022 \\
P3 & dir & p3->p4 & 0.584±0.018 & 0.591±0.029 \\
P3 & all & all->p4 & 0.686±0.114 & 0.659±0.073 \\
P3MS & dir & p12->p4 & 0.655±0.016 & 0.593±0.016 \\
P3MS & dir & p13->p4 & 0.536±0.009 & 0.533±0.008 \\
P3MS & dir & p23->p4 & 0.555±0.005 & 0.537±0.011 \\
P3MS & all & all->p4 & 0.582±0.064 & 0.554±0.034 \\
\hline
\end{tabular}
\end{table}
% </qpaper-block>

% <qpaper-block id="para.result.validation_controls" question="Q.result.05_validation_controls" claims="" evidence="A_results_paper_layout_2026_03_03_table_p3ms_compact_csv,A_results_paper_layout_2026_03_03_table_to4_compact_csv,A_results_preproc20260626_paper_tables_table_preproc_protocol_condition_csv,A_results_preproc20260626_preproc_ablation_report_md">
These transfer results are interpreted alongside validation controls, not as an isolated negative table. Multi-source training, target-specific transfer directions, and preprocessing variants all test whether the observed gap can be attributed to a simple protocol artifact. They reduce some ambiguities, but they do not change the claim boundary: the present system is supported as personalized lexicon-free spelling, not as calibration-free deployment.

% </qpaper-block>

% <qpaper-block id="para.limit.registration" question="Q.limit.registration" claims="C.boundary" evidence="A_results_slt_p2_20260626_kshot_feasibility_csv,A_results_slt_p2_20260626_experiment_status_md,A_results_slt_p2_20260626_split_audit_protocol_split_summary_csv,A_results_preproc20260626_paper_tables_table_preproc_protocol_condition_csv,A_results_preproc20260626_experiment_status_md">
The strongest limitation is not decoding speed or the existence of a recognizer. It is calibration. The experiments show that EPGSpeller can work well when participant-specific data are available, but cross-user transfer remains weak. The preprocessing ablation rules out a simple explanation in which train-only standardization, contact-rate matching, target marginal normalization, or PCA alone closes the gap; the best protocol-level preprocessing averages remain 0.6426 CER for P3 and 0.5692 CER for P3MS. The experiments also show that the present public recordings do not contain enough repeated productions per word to run a standard k=4 or k=8 50-word calibration curve. Future AAC-oriented EPG datasets therefore need denser repeated spelling trials, not only more model variants.

% </qpaper-block>
